\chapter{Contribuciones (taxonom\'ia CRediT)}
\label{cap:credit}

La siguiente tabla asigna categor\'ias de la taxonom\'ia CRediT
(\emph{Contributor Roles Taxonomy})~\cite{credit-taxonomy} a cada
integrante del equipo, \'unicamente sobre el trabajo efectivamente
documentado en el repositorio (c\'odigo, ADR, evidencia de mediciones,
Postman, requisitos). Ninguna categor\'ia se asigna sin respaldo
verificable.

\begin{table}[H]
  \centering
  \small
  \caption{Contribuciones CRediT del equipo BMT en la Tercera Entrega de BIOPET.}
  \label{tab:credit}
  \begin{tabular}{@{}p{4.2cm}ccc@{}}
    \toprule
    Categor\'ia CRediT & Jaime & Fred & Zaida \\
    \midrule
    Conceptualization & \checkmark & \checkmark & \checkmark \\
    Methodology & \checkmark & \checkmark & \checkmark \\
    Software & \checkmark & \checkmark & \checkmark \\
    Validation & \checkmark & \checkmark & \\
    Formal analysis & & \checkmark & \\
    Investigation & \checkmark & \checkmark & \checkmark \\
    Data curation & \checkmark & \checkmark & \\
    Writing -- Original Draft & \checkmark & & \checkmark \\
    Writing -- Review \& Editing & \checkmark & \checkmark & \checkmark \\
    Visualization & \checkmark & \checkmark & \\
    \bottomrule
  \end{tabular}
\end{table}

\section{Detalle por integrante}

\subsection*{Jaime Josu\'e Mariscal Cabrera}
Seguridad del backend (OWASP A01, A02, A03, A05, A07, A09); dise\~no e
implementaci\'on de la autenticaci\'on por cookies, revocaci\'on de tokens,
rate limiting de login y auditor\'ia estructurada; formato uniforme de
error \texttt{ProblemDetail}; pruebas automatizadas asociadas
(\texttt{AuthControllerTest}, \texttt{JwtCookieAuthenticationTest},
\texttt{SqlInjectionSecurityTest}, \texttt{SecurityHeadersTest},
\texttt{AuthenticationAuditServiceTest}, entre otras); integraci\'on y
umbral autom\'atico de cobertura JaCoCo; colecci\'on y entorno Postman
reproducibles; diagrama C4 Nivel~3 de componentes del backend; ADR-006; y
la documentaci\'on de seguridad de \texttt{docs/mediciones/sec/} y de las
secciones de Seguridad y Cobertura JaCoCo del diccionario de datos.

\subsection*{Fred Adri\'an Beltr\'an Montiel}
Estrategia de base de datos reproducible con PostgreSQL (ADR-004):
separaci\'on de privilegios (\texttt{biopet_app}), scripts de
inicializaci\'on de Docker (\texttt{db/schema.sql}, \texttt{db/roles.sql},
\texttt{db/seed.sql}); reproducibilidad del despliegue (ADR-005):
\texttt{Makefile} y pinning de im\'agenes por digest \texttt{sha256};
funci\'on nativa \texttt{fn_resumen_mascotas_por_especie} y su prueba de
integraci\'on con Testcontainers; cach\'e Redis del listado de mascotas;
dise\~no, ejecuci\'on y an\'alisis estad\'istico de las corridas de
rendimiento con k6 (intervalos de confianza, percentiles, tasa de error,
throughput, hit ratio).

\subsection*{Zaida Melissa Taipe Mora}
Consolidaci\'on de la ingenier\'ia de requisitos (SRS, historias de
usuario, casos de uso, \texttt{CHANGELOG-REQ.md}); desarrollo del frontend
Angular (autenticaci\'on por cookies, guardas de ruta, interceptor de
errores); CRUD de mascotas accesible en la interfaz (atributos ARIA,
manejo de foco y anuncios de error para lectores de pantalla); dise\~no
metodol\'ogico de las mediciones de usabilidad (SUS) y accesibilidad
autom\'atica (Lighthouse), pendientes de ejecuci\'on; documentaci\'on
funcional del sistema.

\bigskip
\noindent\textit{Nota:} las tareas de base compartidas de la Entrega~1B
(estructura inicial del CRUD de mascotas en el backend, esquema JWT
original) no se atribuyen individualmente en esta tabla por no existir en
el repositorio de la Tercera Entrega un registro de autor\'ia expl\'icito
por commit que permita separarlas con certeza; se documentan como trabajo
de equipo en la matriz de trazabilidad (cap\'itulo~\ref{cap:trazabilidad}).
